\section*{Chapter \ref{ch:ll:lock-free-programming} --- Lock-Free Programming}

\begin{solution}{exo:ll:lock-free-programming:1}
On the committed measurement about \qty{4.7}{\nano\second} alone and \qty{95}{\nano\second} with four threads; together the four
threads achieve $4/\qty{95}{\nano\second} \approx 42$ million shared increments a second, against about 730 million with
per-thread counters.
\end{solution}

\begin{solution}{exo:ll:lock-free-programming:2}
With release on the flag and acquire on its load, no: the acquire that sees the flag sees every write before the release,
including the price. With a relaxed flag, yes: nothing orders the price's store before the flag's, for the compiler or,
on other architectures, the processor.
\end{solution}

\begin{solution}{exo:ll:lock-free-programming:3}
$2^{32}/10^8 \approx 43$ seconds. A wrap causes a false match only if a thread stalls between its read and its CAS for exactly
a multiple of $2^{32}$ updates and finds the same index at the same tag: possible in theory, vanishingly rare in practice,
and excluded by a 64-bit tag with a 16-byte CAS if needed.
\end{solution}

\begin{solution}{exo:ll:lock-free-programming:4}
Under \omterm{def:ll:lock-free-programming:ordering}{sequential consistency} $(0,1)$, $(1,0)$ and $(1,1)$; x86 with weaker orderings also allows $(0,0)$, when both stores are
still in their cores' store buffers as the loads complete. The window is a few nanoseconds wide, while the two threads' starts
differ by far more on each trial, so they rarely overlap: tens to hundreds of times per million here.
\end{solution}

\begin{solution}{exo:ll:lock-free-programming:5}
A producer claims a position with its \omterm{def:ll:lock-free-programming:cas}{compare-and-swap}, then writes the cell and publishes it. If it is preempted between the
two, the cell's sequence number stays unpublished: the consumer of that position finds the queue ``empty'' and every later
position waits behind it, so one stalled thread stops the others' progress.
\end{solution}

\begin{solution}{exo:ll:lock-free-programming:6}
The mutex is held 20\% of the time ($100\,000 \times \qty{2}{\micro\second}$), so about one order in five waits, a microsecond
on average. If the logger is preempted while holding it, an order can wait the whole quantum, \qty{4}{\milli\second}.
\end{solution}

\begin{solution}{exo:ll:lock-free-programming:7}
In twenty runs of two million operations no node was lost or duplicated (\texttt{measured\_aba.csv}): the vulnerable window,
between reading the successor and the \omterm{def:ll:lock-free-programming:cas}{compare-and-swap}, lasts a few nanoseconds and the scheduler almost never preempts a
thread inside it. The code is still wrong; the deterministic replay of \cref{fig:ll:lock-free-programming:aba} shows the
interleaving. Running \omterm{def:ll:lock-free-programming:progress}{lock-free} code proves little; the argument must come from reasoning, model checking or tools that
explore interleavings.
\end{solution}

\begin{solution}{exo:ll:lock-free-programming:8}
Correctness under the C++ model is not a property of one processor: x86 happens to keep stores in order, but the compiler may
reorder the snapshot's writes after a relaxed flag store, and another architecture or compiler version may. The absence of
failure in a test is not evidence. Use release on the flag and acquire on its load, or a sequence lock (chapter~12).
\end{solution}

\begin{solution}{pb:ll:lock-free-programming:1}
\begin{enumerate}
\item $50\,000 \times \qty{3}{\micro\second} = 15\%$ of the time.
\item 15\% of orders (arrivals see the time average), waiting half a hold on average: \qty{1.5}{\micro\second}.
\item $20\,000 \times 0.15 \times \qty{1.5}{\micro\second} = \qty{4.5}{\milli\second}$ a second.
\item With two threads taking it often, each acquisition costs several times its uncontended cost (the lock's line bounces
between cores): about \qty{100}{\nano\second} against \qty{14}{\nano\second} on the committed measurement.
\item $50\,000 \times 0.0001 = 5$ a second.
\item $5 \times \qty{3}{\milli\second} = \qty{15}{\milli\second}$ a second: 1.5\% of the time.
\item $20\,000 \times 0.015 = 300$ orders a second, each waiting up to \qty{3}{\milli\second}, \qty{1.5}{\milli\second} on
average.
\item $300 \times \qty{1.5}{\milli\second} + \qty{4.5}{\milli\second} \approx \qty{455}{\milli\second}$ a second: the rare
preemptions dominate.
\item It writes the order into its own ring slot and publishes it with a release store: tens of nanoseconds, never waiting for
the logger.
\item The strategy must decide: drop the log record (and count the drop) or overwrite, never block; the ring is sized so that it
does not fill.
\item Wait-freedom: every push completes in a bounded number of its own steps.
\item Priority inheritance raises the logger's priority while it holds the lock, but a trading thread's deadline is microseconds;
waking and running the holder takes longer than that, and preemption by the kernel itself is not a matter of priorities.
\item \textbf{Named result.} With the mutex the strategy is blocked about \qty{455}{\milli\second} a second on average, almost all
of it behind a preempted logger; with the ring, never.
\item They enabled priority inheritance on the mutex by uploading a patch: enough because their deadline was a bus cycle of
milliseconds and the holder could finish quickly once raised.
\item In the standard library (allocators, logging, locale), in third-party libraries, in the dynamic loader, and in any shared
configuration object.
\item A way to know that no thread still reads a node before it is reused: tags on reused nodes, \omterm{def:ll:lock-free-programming:reclaim}{hazard pointers} or epochs on
freed ones.
\item Two producers and two consumers move 300\,000 items and check each arrives once, in C++ and Rust, and the C++ test runs under
the thread sanitiser. They cannot prove the absence of rare interleavings.
\item Off the hot path, for state touched rarely by threads that can afford to wait (configuration, reporting), where simplicity
wins.
\item Release on the store that publishes the slot and acquire on the load that consumes it.
\item It removes waiting on other threads from the hot path: no thread's stall can stop it.
\end{enumerate}
\end{solution}

\begin{solution}{iq:ll:lock-free-programming:1}
Relaxed: atomicity only. Acquire--release: a release store publishes all earlier writes to the thread whose acquire load reads
it. Sequentially consistent: additionally a single total order of all such operations, forbidding for example both threads of a
store-buffer test reading zero.
\iqlookfor{publication versus a global order, with an example.}
\end{solution}

\begin{solution}{iq:ll:lock-free-programming:2}
A \omterm{def:ll:lock-free-programming:cas}{compare-and-swap} succeeds because the value is again what it read, although it changed in between and the state computed from
it is stale. Prevent it with a tag incremented on every update (compared with the pointer), or by never reusing a node while a
reader may hold it (\omterm{def:ll:lock-free-programming:reclaim}{hazard pointers}, epochs, or nodes that are never freed).
\iqlookfor{the mechanism and two independent remedies.}
\end{solution}

\begin{solution}{iq:ll:lock-free-programming:3}
No. Under heavy contention on one location, a CAS loop retries and can be slower than a mutex; the chapter's counters show a CAS
loop at about \qty{200}{\nano\second} per increment with four threads. Lock-freedom is a guarantee about stalls, not about speed.
\iqlookfor{progress versus throughput.}
\end{solution}

\begin{solution}{iq:ll:lock-free-programming:4}
A high-priority thread waiting for a lock held by a low-priority thread that medium-priority work keeps from running. On a desk the
hot thread waits for a logger or reporter that the scheduler preempted, for milliseconds.
\iqlookfor{the three-thread mechanism and a trading instance.}
\end{solution}

\begin{solution}{iq:ll:lock-free-programming:5}
Stress tests with checks of invariants (every item once), sanitisers, tests that perturb scheduling, model checking of the
algorithm, and deterministic replays of known bad interleavings. A stress test cannot show that a rare interleaving is
harmless, only that it did not happen.
\iqlookfor{tools beyond running it, and the limits of testing.}
\end{solution}

\begin{solution}{iq:ll:lock-free-programming:6}
One counter per thread, each on its own \omterm{def:ll:memory-hierarchy-and-caches:line}{cache line}, incremented with a relaxed operation (or a plain write if only its owner
writes and the reader tolerates a slightly stale value); the reader sums them. No line is shared between writers.
\iqlookfor{sharding by thread and cache-line padding.}
\end{solution}
